%==============================================================
% LAMPIRAN
%==============================================================

% Halaman standalone judul "LAMPIRAN" di tengah halaman
\clearpage
\vspace*{\fill}
\begin{center}
  {\fontsize{14pt}{17pt}\selectfont\bfseries LAMPIRAN}
\end{center}
\vspace*{\fill}
\addcontentsline{toc}{chapter}{LAMPIRAN}

%--------------------------------------------------------------
% LAMPIRAN 1: Konfigurasi Hyperparameter Model
%--------------------------------------------------------------
\clearpage

\refstepcounter{lampiran}
\addcontentsline{lamp}{lampiran}{\protect\numberline{\thelampiran}Lampiran 1\enspace
Konfigurasi \textit{Hyperparameter} Model TinyResNet}
\label{lamp:hyperparameter}

\noindent\hangindent=5.6em
\makebox[5.4em][l]{\textbf{Lampiran \thelampiran}}%
Konfigurasi \textit{Hyperparameter} Model TinyResNet pada Eksperimen Terbaik

\vspace{4pt}

\centering
\begin{tabular}{l l}
  \toprule
  \textbf{Parameter} & \textbf{Nilai} \\
  \midrule
  Arsitektur          & TinyResNet (3 stage residual) \\
  Lebar kanal         & 48 / 96 / 192 \\
  Ukuran input        & $224 \times 224$ piksel \\
  Jumlah kelas        & 5 \\
  Fungsi kerugian     & Focal Loss ($\gamma = 1{,}0$) \\
  Bobot kelas         & Berbanding terbalik dengan frekuensi \\
  Optimizer           & AdamW \\
  Laju belajar        & $10^{-3}$ \\
  Peluruhan bobot     & $5 \times 10^{-4}$ \\
  Penjadwalan LR      & \textit{Cosine annealing} \\
  Jumlah epoch        & 160 \\
  Ukuran batch        & 32 \\
  Augmentasi MixUp    & $\alpha = 0{,}2$, $p = 0{,}5$ \\
  Augmentasi CutMix   & $\alpha = 1{,}0$, $p = 0{,}5$ \\
  \textit{Label smoothing} & 0,1 \\
  \textit{Stochastic depth} & $0 \rightarrow 0{,}2$ \\
  Weight EMA decay    & 0,999 \\
  Protokol CV         & StratifiedGroupKFold, 5 \textit{fold} \\
  Platform komputasi  & Modal (A100-80GB) \\
  \bottomrule
\end{tabular}

\clearpage

%--------------------------------------------------------------
% LAMPIRAN 2: Konfigurasi Representasi Softmap Spasial
%--------------------------------------------------------------

\refstepcounter{lampiran}
\addcontentsline{lamp}{lampiran}{\protect\numberline{\thelampiran}Lampiran 2\enspace
Konfigurasi Representasi Softmap Spasial}
\label{lamp:vessel_config}

\noindent\hangindent=5.6em
\makebox[5.4em][l]{\textbf{Lampiran \thelampiran}}%
Konfigurasi Representasi Softmap Spasial

\vspace{4pt}

\noindent
\begin{tabularx}{\textwidth}{l X}
  \toprule
  \textbf{Parameter} & \textbf{Nilai} \\
  \midrule
  Ukuran kerja maksimum & Sisi terpanjang 768 piksel \\
  Ukuran input model & $224 \times 224$ piksel \\
  Masking FOV & Ambang intensitas rendah, closing, fill holes, komponen terbesar \\
  Kanal 1 & \textit{Plain-Gabor vessel softmap} \\
  Vessel CLAHE awal & \textit{Clip limit} 6,0; \textit{tile} $16 \times 16$ \\
  Gabor: sigma & 1,5; 2,5; 3,5; 5,0 piksel \\
  Gabor: lambda & 3; 5; 7; 10 piksel \\
  Gabor: sudut & 12 sudut ($0^\circ$--$165^\circ$, langkah $15^\circ$) \\
  Vessel median blur & Kernel $7 \times 7$ \\
  Vessel CLAHE akhir & \textit{Clip limit} 12,0; \textit{tile} $12 \times 12$ \\
  Kanal 2 & \textit{Ridge softmap} \\
  Meijering ridge: sigma & 3; 5; 7; 9; 11 piksel \\
  Oriented top-hat & 12 orientasi, elemen struktur garis \\
  Formula ridge & $0{,}60 \times \text{meijering\_ridge} + 0{,}40 \times \text{oriented\_tophat}$ \\
  Kanal 3 & Kanal hijau CLAHE yang dimasker FOV \\
  Normalisasi kanal & Persentil 1--99 ke rentang $[0, 1]$ \\
  Evaluasi pendukung segmentasi & Dice pembuluh darah 0,4739 pada HVDROPDB-BV; Dice \textit{ridge} pada HVDROPDB-RIDGE \\
  \bottomrule
\end{tabularx}

\clearpage

%--------------------------------------------------------------
% LAMPIRAN 3: Distribusi Dataset Zhao2024
%--------------------------------------------------------------

\refstepcounter{lampiran}
\addcontentsline{lamp}{lampiran}{\protect\numberline{\thelampiran}Lampiran 3\enspace
Distribusi Lengkap Dataset Zhao2024}
\label{lamp:dataset}

\noindent\hangindent=5.6em
\makebox[5.4em][l]{\textbf{Lampiran \thelampiran}}%
Distribusi Lengkap Dataset Zhao2024 dan Statistik Kelompok

\vspace{4pt}

\noindent
\begin{tabular}{l c c c}
  \toprule
  \textbf{Kelas} & \textbf{Jumlah Citra} & \textbf{Persentase} & \textbf{Jumlah Kelompok} \\
  \midrule
  Laser Scars & 343 & 31,2\%  & 343 (singleton) \\
  Stage 3     & 261 & 23,7\%  & $\sim$120 \\
  Normal      & 236 & 21,5\%  & $\sim$100 \\
  Stage 2     & 165 & 15,0\%  & $\sim$80 \\
  Stage 1     & 94  &  8,6\%  & $\sim$79 \\
  \midrule
  \textbf{Total} & \textbf{1.099} & \textbf{100\%} & \textbf{$\sim$721} \\
  \bottomrule
\end{tabular}

\vspace{6pt}
\noindent Catatan: jumlah kelompok merupakan perkiraan berdasarkan analisis NCC $\geq$ 0,99.
Bekas luka laser (\textit{laser scars}) diperlakukan sebagai \textit{singleton group} karena
tidak memiliki kecocokan di kelas lain.
